\section{Related Work}\label{sec:related}

\paragraph{How language models answer multiple-choice questions.}
Circuit analyses find that attention heads select the answer among the listed options.
\citet{lieberum2023chinchilla} identify ``correct letter'' heads in Chinchilla-70B whose query and key subspaces partly encode an ``$N$th item in an enumeration'' feature; \citet{wiegreffe2025answer} attribute answer-symbol prediction to the attention of a few middle layers; \citet{tulchinskii2024wise} find middle-layer select-and-copy heads whose pre-RoPE query--key score between the last token and the end of each option can recover the correct option more accurately than the model's own answer; and \citet{kamath2025attention} explain attention patterns through query--key feature interactions.
\citet{sun2026persuaded} trace persuasion to decision heads that copy the option they attend to, steered by a routing feature in the option tokens, and \citet{wong2026decide} describe a two-stage computation that selects a winner in content space before binding it to an output symbol.
\citet{ok2026lost} show that when options precede the context, the causal mask prevents option tokens from attending to it, the constraint our \fmtBefore{} control exploits.
These works mostly study how the answer is chosen among options; we study how option tokens read a state written earlier in the context, and show that their presence after the state token routes part of the state, and from 3B on most of it, through its key rather than its value (\cref{tab:natural}).

\paragraph{Binding and in-context state retrieval.}
Models bind entities to attributes with binding ID vectors \citep{feng2023binding}, and track beliefs with lookbacks in which ordering IDs at the state token let a binding lookback retrieve the state's ID and an answer lookback retrieve the state token \citep{prakash2025lookbacks}.
\citet{oh2026rebinding} find a rebinding circuit whose binding signal lies in query/key subspaces or mainly in key vectors, depending on the model family; \citet{wu2026record} find that the model reads mainly the address, not the value, recorded in an operation's KV cache; and \citet{steele2026routing,steele2026spaces} separate a shared value slot from a router that selects which frame (a character's belief or reality) or, more generally, which mental space is read.
\citet{bojieli2026notes} show that overwriting one field's cached keys and values while reusing the downstream cache barely changes the decision, because later tokens have already written their conclusions at prefill.
We vary only the readout and find that the state is copied mainly from the state token's value when nothing re-mentions the candidates (key identity at most \noneMax{} nats), but is also looked up through its key when they are listed as options after it, in all \nModels{} models, and, from 7B on, when a neutral sentence re-mentions them. This identity-matching route, whose readers at 1.5B are the option words themselves (\cref{sec:res-readers}), sits alongside the address and payload accounts above.

\paragraph{Validity of interventions and causal abstraction.}
Interchange interventions are used to train a network to realise a given alignment between causal variables and representations \citep{geiger2021inducing} and to learn such an alignment in a distributed subspace \citep{geiger2024finding}.
Whether such an intervention uses the model's own mechanism is debated: subspace patches can act through dormant parallel pathways \citep{makelov2024subspace}, a finding \citet{wu2024reply} argue is an artifact of the training and evaluation paradigm; optimised interventions can reach the right answer through planted pathways \citep{liu2026wrong}; successful steering need not reproduce the natural activation \citep{shih2026steering}; and steering vectors are not identifiable from behaviour \citep{venkatesh2026steering}.
Diagnostics and benchmarks include control tasks \citep{hewitt2019control}, RAVEL \citep{huang2024ravel}, CausalGym \citep{arora2024causalgym}, MIB \citep{mueller2025mib} and input-space partitions that show where an abstraction is faithful \citep{li2026bucketing}.
\citet{anonymous2026fitted} fit rank-16 interventions for source transfer and for a location remap on belief stories, and find that exchanging the critical-token attention keys induced by learned and PCA patches transfers a substantial part of their behavioural difference in both directions.
These works ask whether an intervention is faithful or identifiable, and \citet{anonymous2026fitted} attribute part of one intervention's effect to keys under their own readout; we show that for that remap the part carried by critical-token keys depends on the readout while the behavioural effect barely does ($\varphi$ from \phiOptMistral{} to \phiLetterMistral{} at Mistral-Small-24B and from \phiNoneQwen{} to \phiBeforeQwen{} at Qwen2.5-72B; \cref{tab:frames}).

\paragraph{Format dependence of representations and steering.}
Performance is sensitive to meaning-preserving changes in prompt format \citep{sclar2023quantifying}.
Function vectors extracted from open-ended and multiple-choice formats are nearly orthogonal \citep{opielka2026causality}; steering directions can follow answer identifiers rather than the intended judgment, and multiple-choice and open-ended evaluations can yield different behavioural conclusions \citep{gao2026encodings}.
At the component level, refusal steering interacts with attention mainly through the OV circuit rather than the QK circuit and loses little of its effect when attention scores are frozen \citep{cheng2026steering}, and more accurate introspective localisation of an injected concept vector is associated with better alignment of the key and value changes it induces with the QK and OV computations of gate heads \citep{zou2026introspection}.
These works show that behaviour or representations differ across formats, or relate an intervention's effect to QK and OV computations; we fix one intervention whose behavioural effect is nearly format-invariant and show that the part of it carried by critical-token keys is not ($\psi$ from \psiLetterQwen{} to \psiBeforeQwen{} at Qwen2.5-72B), so a QK-versus-OV attribution of an intervention can depend on the readout under which it is measured.

\paragraph{Activation patching and KV-cache edits.}
Component-level attribution rests on causal mediation analysis \citep{vig2020causal} and intervention-based circuit discovery \citep{wang2022ioi}, whose conclusions depend on choices of metric and corruption \citep{zhang2023patching,heimersheim2024patching}.
Precomputed KV caches are exact substitutes for a joint forward pass because of the causal mask, and contrastive edits to cached values can steer behaviour whereas key arithmetic destroys coherence \citep{pustovit2026packs}.
We use the same causal-mask argument to make the joint key--value clamp exact, but edit the key and the value of one context token separately, splice attention by row so that only chosen later tokens see a swapped key (\cref{sec:method-splice}), and add the readout to the patching choices that can change a component-level conclusion.
